\documentclass[11pt,a4paper]{article}

% ─── Packages ────────────────────────────────────────────────────────────────
\usepackage[T1]{fontenc}
\usepackage[utf8]{inputenc}
\usepackage[english]{babel}
\usepackage{lmodern}
\usepackage{microtype}
\usepackage{geometry}
\geometry{margin=2.2cm, top=2.5cm, bottom=2.5cm}

\usepackage{graphicx}
\usepackage{float}
\usepackage{booktabs}
\usepackage{longtable}
\usepackage{array}
\usepackage{multirow}
\usepackage{xcolor}
\usepackage{colortbl}
\usepackage{hyperref}
\usepackage{enumitem}
\usepackage{parskip}
\usepackage{titlesec}
\usepackage{fancyhdr}
\usepackage{tocloft}
\usepackage{amsmath}
\usepackage{listings}
\usepackage{caption}
\usepackage{subcaption}
\usepackage{mdframed}
\usepackage{tikz}  % for diagram placeholders

% ─── Colors ──────────────────────────────────────────────────────────────────
\definecolor{sfblue}{RGB}{0, 112, 210}
\definecolor{sfnavy}{RGB}{22, 50, 92}
\definecolor{sfgray}{RGB}{112, 112, 112}
\definecolor{sflightgray}{RGB}{243, 243, 243}
\definecolor{sfgreen}{RGB}{45, 160, 75}
\definecolor{sforange}{RGB}{255, 140, 0}
\definecolor{sfred}{RGB}{194, 57, 52}

% ─── Hyperlinks ───────────────────────────────────────────────────────────────
\hypersetup{
    colorlinks=true,
    linkcolor=sfblue,
    urlcolor=sfblue,
    citecolor=sfnavy,
    pdftitle={CLV Prediction Platform — Technical Report},
    pdfauthor={CLV Project Team},
}

% ─── Section Formatting ───────────────────────────────────────────────────────
\titleformat{\section}{\large\bfseries\color{sfnavy}}{{\thesection}}{0.8em}{}[\vspace{-4pt}\rule{\linewidth}{0.6pt}]
\titleformat{\subsection}{\normalsize\bfseries\color{sfblue}}{\thesubsection}{0.6em}{}
\titleformat{\subsubsection}{\normalsize\itshape\color{sfgray}}{\thesubsubsection}{0.5em}{}

% ─── Header / Footer ──────────────────────────────────────────────────────────
\pagestyle{fancy}
\fancyhf{}
\fancyhead[L]{\small\color{sfgray} CLV Prediction Platform — Technical Report}
\fancyhead[R]{\small\color{sfgray} \today}
\fancyfoot[C]{\small\color{sfgray}\thepage}
\renewcommand{\headrulewidth}{0.3pt}

% ─── Code Listings ────────────────────────────────────────────────────────────
\lstset{
    basicstyle=\small\ttfamily,
    backgroundcolor=\color{sflightgray},
    frame=single,
    rulecolor=\color{sfgray},
    breaklines=true,
    columns=flexible,
    keepspaces=true,
    showstringspaces=false,
    keywordstyle=\color{sfblue}\bfseries,
    commentstyle=\color{sfgray}\itshape,
    stringstyle=\color{sfgreen},
}

% ─── Callout Box ──────────────────────────────────────────────────────────────
\newmdenv[
  topline=false, bottomline=false, rightline=false,
  linewidth=3pt, linecolor=sfblue,
  backgroundcolor=sflightgray,
  innerleftmargin=8pt, innerrightmargin=8pt,
  innertopmargin=6pt, innerbottommargin=6pt,
]{callout}

% ─── Placeholder command ──────────────────────────────────────────────────────
\newcommand{\diagramplaceholder}[2]{%
  \begin{center}
    \begin{tikzpicture}
      \fill[sflightgray] (0,0) rectangle (#1,#2);
      \draw[sfgray, dashed, thick] (0,0) rectangle (#1,#2);
      \node[align=center, text=sfgray] at ({#1/2},{#2/2})
        {\textit{[Diagram placeholder]}};
    \end{tikzpicture}
  \end{center}
}

% ─────────────────────────────────────────────────────────────────────────────
\begin{document}

% ─── Title Page ───────────────────────────────────────────────────────────────
\begin{titlepage}
  \centering
  \vspace*{2cm}
  {\color{sfnavy}\rule{\linewidth}{1.2pt}}\\[0.6em]
  {\Huge\bfseries\color{sfnavy} CLV Prediction Platform}\\[0.4em]
  {\Large\color{sfblue} Technical Report}\\[0.3em]
  {\color{sfnavy}\rule{\linewidth}{1.2pt}}\\[1.2em]
  {\large Customer Lifetime Value \& Churn Analytics}\\[0.2em]
  {\large Salesforce CRM-Integrated Web Dashboard}\\[2cm]
  \begin{tabular}{rl}
    \textbf{Version:}   & 1.0 \\[3pt]
    \textbf{Date:}      & \today \\[3pt]
    \textbf{Status:}    & Final Draft \\[3pt]
    \textbf{Scope:}     & 1{,}200 Accounts · B2C Retail Portfolio \\[3pt]
    \textbf{Currency:}  & Czech Koruna (Kč) \\
  \end{tabular}
  \vfill
  {\small\color{sfgray} Confidential — Internal Use Only}
\end{titlepage}

% ─── Table of Contents ────────────────────────────────────────────────────────
\tableofcontents
\newpage

% ═══════════════════════════════════════════════════════════════════════════════
\section{Executive Summary}
% ═══════════════════════════════════════════════════════════════════════════════

This report documents the architecture, technology stack, data lineage, user testing
outcomes, design trade-offs, and deployment recommendations for the
\textbf{CLV Prediction Platform}—an end-to-end machine learning pipeline and
interactive web dashboard built to support B2C retail customer relationship
management.

The system ingests historical Salesforce CRM data (2022--2024), engineers 29
analytical features per account, trains an XGBoost regression model and an
XGBoost churn classifier, and exposes predictions through a Flask REST API.
A single-page application frontend renders the results in a Salesforce Lightning–
inspired dashboard used by analysts, retention specialists, and executives.

\begin{callout}
\textbf{Key Metrics (Production Model, n = 1{,}200 accounts):}\\[4pt]
\begin{tabular}{ll}
  MAE (full portfolio):  & 8{,}172 Kč \\
  MAE (active buyers):   & 13{,}888 Kč \\
  R² (full portfolio):   & 0.521 \\
  Churn Classifier AUC:  & 0.883 \\
  Tier Accuracy:         & 48.3\% (current CRM tiers vs.\ model-suggested) \\
  Revenue at Risk:       & 26{,}084 Kč \\
  Upsell Value Potential:& 3{,}765{,}144 Kč \\
\end{tabular}
\end{callout}

% ═══════════════════════════════════════════════════════════════════════════════
\section{System Architecture}
% ═══════════════════════════════════════════════════════════════════════════════

\subsection{Architecture Overview}

The platform follows a three-tier architecture: a \textbf{data processing} layer
(Jupyter notebooks), a \textbf{backend serving} layer (Flask), and a
\textbf{presentation} layer (single-page web application).

\begin{figure}[H]
  \caption{System Architecture Diagram — Placeholder}
  \diagramplaceholder{14}{6}
  \label{fig:architecture}
\end{figure}

\subsection{Component Description}

\begin{description}[leftmargin=1.5cm, labelwidth=2.8cm]
  \item[\textbf{Data Layer}]
    Four raw CSV files exported from Salesforce CRM are consumed by
    Jupyter notebooks (\texttt{krok\_01} through \texttt{krok\_08}).
    Notebooks produce cleaned feature files (\texttt{step\_02\_features.csv})
    and churn scores (\texttt{step\_08\_churn\_output.csv}) written to
    \texttt{notebooks/outputs/}.

  \item[\textbf{Model Layer}]
    Serialised XGBoost models (\texttt{xgboost\_clv\_model.pkl},
    \texttt{feature\_names.pkl}) are persisted in \texttt{models/}.
    On first request, \texttt{app.py} trains the model in memory; subsequent
    requests load from disk, reducing latency to milliseconds.

  \item[\textbf{API Layer}]
    A Flask application (\texttt{app.py}) exposes five REST endpoints
    consumed by the frontend via \texttt{fetch()}.

  \item[\textbf{Frontend Layer}]
    A single HTML page (\texttt{templates/index.html}) with vanilla
    JavaScript (\texttt{static/script.js}) and CSS (\texttt{static/style.css}).
    Chart.js renders all eight analytical charts.
\end{description}

\subsection{API Endpoints}

\begin{table}[H]
\centering
\small
\begin{tabular}{lp{3.5cm}p{6cm}}
\toprule
\textbf{Method} & \textbf{Endpoint} & \textbf{Description} \\
\midrule
GET  & \texttt{/api/accounts}      & Full account list with features, churn probability, and CLV prediction \\
POST & \texttt{/api/predict}        & Triggers model training or loading; runs predictions for all accounts \\
POST & \texttt{/api/simulate}       & Co-pilot scenario simulation — overrides selected metrics and re-predicts \\
GET  & \texttt{/api/explain/<id>}   & XAI endpoint — returns top-3 positive/negative CLV drivers for one account \\
GET  & \texttt{/api/segments}       & Computes and returns 5 RFM-based micro-segment cohorts \\
\bottomrule
\end{tabular}
\caption{REST API surface}
\end{table}

% ═══════════════════════════════════════════════════════════════════════════════
\section{Technology Stack \& Data Lineage}
% ═══════════════════════════════════════════════════════════════════════════════

\subsection{Technology Stack}

\begin{table}[H]
\centering
\small
\begin{tabular}{lll}
\toprule
\textbf{Layer} & \textbf{Technology} & \textbf{Version / Note} \\
\midrule
\multirow{6}{*}{Backend / ML} 
  & Python           & 3.11+ \\
  & Flask            & $\geq$ 2.3.0 \\
  & Flask-CORS       & $\geq$ 4.0.0 \\
  & XGBoost          & $\geq$ 1.7.0 \\
  & scikit-learn     & $\geq$ 1.3.0 \\
  & pandas / numpy   & $\geq$ 2.0.0 / 1.24.0 \\
\midrule
\multirow{3}{*}{Analysis}
  & SHAP             & $\geq$ 0.42.0 \\
  & matplotlib / seaborn & $\geq$ 3.7.0 / 0.12.0 \\
  & scipy            & $\geq$ 1.11.0 \\
\midrule
\multirow{2}{*}{Persistence}
  & joblib           & $\geq$ 1.3.0 \\
  & openpyxl         & $\geq$ 3.1.0 \\
\midrule
\multirow{2}{*}{Frontend}
  & Chart.js         & v4 (CDN) \\
  & Vanilla JS / CSS & ES2020, CSS custom properties \\
\midrule
Deployment & gunicorn & Latest stable \\
\bottomrule
\end{tabular}
\caption{Full technology stack}
\end{table}

\subsection{Data Lineage}

\begin{figure}[H]
  \caption{Data Lineage Diagram — Placeholder}
  \diagramplaceholder{14}{7}
  \label{fig:lineage}
\end{figure}

The data flows through the following stages:

\begin{enumerate}[leftmargin=1.5cm]
  \item \textbf{Raw Ingestion} — Four Salesforce CSV exports are placed in \texttt{csv/}:
    \begin{itemize}[noitemsep]
      \item \texttt{Account.csv} — 1{,}200 customer entities (demographics, tier, channel)
      \item \texttt{Order\_\_c.csv} — $\sim$9{,}500 transaction records (2022--2024)
      \item \texttt{Activity\_\_c.csv} — 1{,}200 digital engagement records
      \item \texttt{Product2.csv} — 72 product catalogue records
    \end{itemize}

  \item \textbf{EDA \& Cleaning} (\texttt{krok\_01\_EDA.ipynb}) — Null handling,
    outlier analysis, distribution checks; outputs \texttt{step\_01\_*.csv}.

  \item \textbf{Feature Engineering} (\texttt{krok\_02\_feature\_engineering.ipynb}) —
    RFM aggregation, spend trends, engagement normalisation, category diversity.
    Output: \texttt{step\_02\_features.csv} (1{,}200 rows $\times$ 33 columns).

  \item \textbf{Model Training}:
    \begin{itemize}[noitemsep]
      \item \texttt{krok\_03}: Linear Regression baseline (MAE 9{,}708 Kč, R² 0.418)
      \item \texttt{krok\_04}: Logistic Regression — purchase-probability classifier (AUC 0.910)
      \item \texttt{krok\_05}: Random Forest tuned (MAE 8{,}406 Kč, R² 0.477)
      \item \texttt{krok\_06}: \textbf{XGBoost regressor — final model} (MAE 8{,}172 Kč, R² 0.521)
      \item \texttt{krok\_07}: Customer segmentation logic
      \item \texttt{krok\_08}: \textbf{XGBoost churn classifier} (AUC 0.883, F1 0.746)
    \end{itemize}

  \item \textbf{Runtime Serving} — \texttt{app.py} loads serialised models from
    \texttt{models/}, merges live account data with pre-computed features, and
    applies churn-adjusted CLV:
    \[
      \text{CLV}_{\text{adj}} = \hat{\text{CLV}} \times \bigl(1 - P(\text{churn})\bigr)
    \]

  \item \textbf{Downstream Export} — Users may export predictions as
    \texttt{clv\_predictions\_export.csv} for re-import into Salesforce CRM.
\end{enumerate}

\subsection{Feature Set (29 predictive features)}

\begin{table}[H]
\centering\small
\begin{tabular}{ll}
\toprule
\textbf{Category} & \textbf{Features} \\
\midrule
RFM            & recency\_days, frequency, monetary\_total, monetary\_avg \\
Historical spend & spend\_2022, spend\_2023, spend\_2024 \\
Spend trends   & spend\_trend\_2y, spend\_trend\_1y \\
Digital engagement & login\_count\_30d, login\_count\_90d, email\_open\_rate, \\
                   & app\_usage\_score, days\_since\_login \\
Customer profile & age, tenure\_days, category\_diversity, \\
                 & preferred channel (encoded), campaign subscription status \\
\bottomrule
\end{tabular}
\caption{Feature groups used by the XGBoost CLV regressor}
\end{table}

% ═══════════════════════════════════════════════════════════════════════════════
\section{Model Design}
% ═══════════════════════════════════════════════════════════════════════════════

\subsection{CLV Regressor — XGBoost}

The primary model is an XGBoost gradient boosting regressor trained on
80\% of the 1{,}200-account dataset and evaluated on the held-out 20\%.
Key hyperparameters (found via GridSearch with 5-fold CV):

\begin{lstlisting}[language=Python]
XGBRegressor(
    n_estimators   = 300,
    max_depth       = 6,
    learning_rate   = 0.05,
    random_state    = 42,
)
\end{lstlisting}

\begin{table}[H]
\centering\small
\begin{tabular}{lrrrr}
\toprule
\textbf{Model} & \textbf{MAE (Kč)} & \textbf{RMSE (Kč)} & \textbf{R²} & \textbf{MAE active} \\
\midrule
Linear Regression (baseline) & 9{,}708 & 17{,}101 & 0.418 & 15{,}185 \\
Random Forest (tuned)        & 8{,}406 & 16{,}212 & 0.477 & 14{,}153 \\
\textbf{XGBoost (tuned)}     & \textbf{8{,}172} & \textbf{15{,}512} & \textbf{0.521} & \textbf{13{,}888} \\
\bottomrule
\end{tabular}
\caption{Regression model comparison (test set, 20\% holdout)}
\end{table}

\subsection{Churn Classifier — XGBoost}

A binary churn classifier outputs $P(\text{churn}) \in [0,1]$ per account.
Thresholds applied to the dashboard:

\begin{center}
  Low Risk: $P \leq 35\%$ \quad|\quad
  Medium Risk: $35\% < P \leq 65\%$ \quad|\quad
  High Risk: $P > 65\%$
\end{center}

\begin{table}[H]
\centering\small
\begin{tabular}{lrrrr}
\toprule
\textbf{Model} & \textbf{AUC} & \textbf{F1} & \textbf{Accuracy} & \textbf{Precision} \\
\midrule
Logistic Regression (baseline) & 0.853 & 0.725 & 77.4\% & 76.0\% \\
\textbf{XGBoost (tuned)}       & \textbf{0.883} & \textbf{0.746} & \textbf{79.8\%} & \textbf{80.7\%} \\
\bottomrule
\end{tabular}
\caption{Churn classifier comparison}
\end{table}

\subsection{Loyalty Tier Thresholds}

Thresholds were set using percentile analysis of actual 2025 spending:

\begin{center}
  \textbf{Bronze}: $\hat{\text{CLV}} < 5{,}000$ Kč \quad
  \textbf{Silver}: $5{,}000 \leq \hat{\text{CLV}} < 25{,}000$ Kč \quad
  \textbf{Gold}: $\hat{\text{CLV}} \geq 25{,}000$ Kč
\end{center}

The 25th percentile of active buyers (4{,}938 Kč) was rounded to 5{,}000 Kč;
the 75th percentile (26{,}315 Kč) was rounded to 25{,}000 Kč.

% ═══════════════════════════════════════════════════════════════════════════════
\section{Customer Segmentation}
% ═══════════════════════════════════════════════════════════════════════════════

Five RFM-based micro-segments are computed at runtime via \texttt{/api/segments}.
Segments are non-mutually-exclusive priority cohorts evaluated in the order shown:

\begin{table}[H]
\centering\small
\begin{tabular}{p{3.2cm}p{5.5cm}p{4cm}}
\toprule
\textbf{Segment} & \textbf{Definition} & \textbf{Recommended Action} \\
\midrule
Champions
  & Gold tier · Low churn · frequency $>$ median (3.0) · recency $<$ 180 days
  & Exclusive previews, loyalty lock-in \\
\midrule
High-Value At-Risk
  & Gold/Silver tier · High or Medium churn risk
  & Immediate win-back, high-touch support \\
\midrule
Loyal Mid-Tier
  & Silver/Bronze tier · Low churn · frequency $\geq$ 2
  & Regular campaigns, relationship maintenance \\
\midrule
Dormant High-Potential
  & Monetary total $>$ median · recency $>$ 180 days
  & Reactivation incentives, survey \\
\midrule
Growing New
  & spend\_trend\_1y $>$ 20\% · predicted CLV $>$ 0
  & Nurture, early tier introduction \\
\bottomrule
\end{tabular}
\caption{Customer micro-segment definitions}
\end{table}

% ═══════════════════════════════════════════════════════════════════════════════
\section{User Testing}
% ═══════════════════════════════════════════════════════════════════════════════

\subsection{Testing Approach}

User testing was conducted informally during iterative development, following a
\textbf{think-aloud} methodology with the primary stakeholder (analyst role).
Feedback was collected in real-time conversation and incorporated into the
next build iteration.

\subsection{Test Sessions}

\begin{table}[H]
\centering\small
\begin{tabular}{p{3cm}p{3cm}p{7cm}}
\toprule
\textbf{Session} & \textbf{Role} & \textbf{Focus Area} \\
\midrule
Session 1 & Data Analyst    & Dashboard navigation, Action Center readability \\
Session 2 & Data Analyst    & Segment Explorer cards — layout, metric selection \\
Session 3 & Data Analyst    & Chart aesthetics, color scheme, data interpretation \\
Session 4 & Data Analyst    & Cross-tab navigation from segment pills \\
\bottomrule
\end{tabular}
\caption{User testing sessions}
\end{table}

\subsection{Findings \& Resolutions}

\begin{enumerate}[leftmargin=1.5cm]
  \item \textbf{Missing Segment Explorer} — The Champions segment card was not
    rendering due to \texttt{NaN} values causing JSON serialisation failure in
    \texttt{/api/segments}. \textit{Resolution:} Added server-side NaN
    sanitisation in \texttt{app.py} before \texttt{jsonify()}.

  \item \textbf{Portfolio value not meaningful} — The aggregate portfolio sum
    displayed per segment card was not interpretable at the cohort level.
    \textit{Resolution:} Replaced with a \textbf{Share \%} metric showing each
    segment's proportion of the total 1{,}200-account base.

  \item \textbf{Segment card colour overload} — Colour-coded cards distracted
    from data; the analyst preferred a neutral, professional appearance.
    \textit{Resolution:} Refactored to a Salesforce Lightning–inspired neutral
    gray palette with blue accents.

  \item \textbf{Inconsistent chart colours} — Charts used a mixed palette that
    hindered comparison across charts.
    \textit{Resolution:} Unified all 8 charts to a single blue-gradient theme
    (white $\to$ dark blue) applied consistently in \texttt{createCharts()}.

  \item \textbf{Segment-to-account navigation} — Users expected clicking on a
    top account pill in a segment card to jump to that account.
    \textit{Resolution:} Implemented \texttt{inspectAccount(accountId)} in
    \texttt{script.js} — switches tab, applies filter, and opens detail panel.
\end{enumerate}

\subsection{Outstanding UX Items}

\begin{itemize}[leftmargin=1.5cm, noitemsep]
  \item Mobile-responsive layout not yet tested (desktop-first design)
  \item Accessibility audit (WCAG 2.1 AA) not completed
  \item Executive role — strategic overview tab requires dedicated user testing
\end{itemize}

% ═══════════════════════════════════════════════════════════════════════════════
\section{Design Trade-offs}
% ═══════════════════════════════════════════════════════════════════════════════

\begin{table}[H]
\centering\small
\begin{tabular}{p{3.5cm}p{5cm}p{4.5cm}}
\toprule
\textbf{Decision} & \textbf{Chosen Approach} & \textbf{Alternative / Trade-off} \\
\midrule
Model selection
  & XGBoost regressor (best MAE/R²)
  & Neural network could capture more complex patterns but requires more data and is less interpretable \\
\midrule
XAI method
  & Approximate feature importance (Z-score + XGBoost weights)
  & Full SHAP values are more rigorous but add significant per-request latency ($\sim$200–500ms) \\
\midrule
Frontend framework
  & Vanilla JS + Chart.js
  & React/Vue would improve maintainability at scale; vanilla JS reduces complexity for a single-user tool \\
\midrule
Data persistence
  & In-memory cache (\texttt{data\_cache} dict)
  & A database (SQLite/PostgreSQL) would enable multi-user concurrency but adds infrastructure overhead \\
\midrule
Churn integration
  & Pre-computed CSV (\texttt{step\_08\_churn\_output.csv})
  & Real-time classifier inference at request time is more accurate but adds model-loading latency \\
\midrule
Tier thresholds
  & Percentile-based (5 000 / 25 000 Kč)
  & Revenue-optimised thresholds via grid search could improve tier accuracy but reduce business interpretability \\
\midrule
Deployment server
  & Flask dev server (single-threaded)
  & Gunicorn (included in \texttt{requirements.txt}) for production; dev server is not safe for concurrent users \\
\midrule
Non-buyers in model
  & Included (639 accounts with CLV = 0)
  & Excluding them improves R² on active buyers (0.383 vs 0.521) but risks under-predicting dormant accounts \\
\bottomrule
\end{tabular}
\caption{Key design decisions and trade-offs}
\end{table}

% ═══════════════════════════════════════════════════════════════════════════════
\section{Deployment Recommendations}
% ═══════════════════════════════════════════════════════════════════════════════

\subsection{Minimum Viable Production Setup}

\begin{enumerate}[leftmargin=1.5cm]
  \item \textbf{Application server} — Replace the Flask development server with
    \textbf{Gunicorn} (already in \texttt{requirements.txt}):
\begin{lstlisting}
gunicorn -w 4 -b 0.0.0.0:8000 app:app
\end{lstlisting}
    Use \textbf{nginx} as a reverse proxy with TLS termination.

  \item \textbf{Environment variables} — Move all paths and any secrets to
    environment variables or a \texttt{.env} file; do not hardcode paths.

  \item \textbf{Model versioning} — Tag serialised model files with a timestamp
    or git commit SHA (e.g., \texttt{xgboost\_clv\_v20250101.pkl}) to allow
    rollback.

  \item \textbf{Data refresh schedule} — Re-run notebooks \texttt{krok\_01}
    through \texttt{krok\_08} monthly or after each CRM data export to keep
    predictions current. Automate with a cron job or Airflow DAG.

  \item \textbf{Access control} — The current application has no authentication.
    Wrap behind an SSO provider (e.g., Salesforce OAuth, Azure AD) before
    exposing to broader business users.
\end{enumerate}

\subsection{Recommended Production Architecture}

\begin{figure}[H]
  \caption{Production Architecture Diagram — Placeholder}
  \diagramplaceholder{14}{5}
  \label{fig:prod_arch}
\end{figure}

\begin{description}[leftmargin=1.5cm]
  \item[Tier 1 — CDN / Load Balancer] nginx with TLS; static assets
    (\texttt{style.css}, \texttt{script.js}) served from CDN.
  \item[Tier 2 — Application Server] Gunicorn + Flask on a Linux VM or
    container (Docker recommended); 4 workers for concurrency.
  \item[Tier 3 — Data / Model Store] Models and CSVs on a shared volume
    or object storage (e.g., Azure Blob, AWS S3). Long-term: migrate to a
    relational database (PostgreSQL) for account data.
  \item[Tier 4 — Batch Pipeline] Jupyter notebooks orchestrated via
    Airflow or Azure Data Factory, triggered monthly on new CRM exports.
\end{description}

\subsection{Scaling Considerations}

\begin{itemize}[leftmargin=1.5cm, noitemsep]
  \item \textbf{Read performance}: The current in-memory \texttt{data\_cache}
    dict works well for a single-user dashboard ($<$ 5 concurrent users). For
    more users, move to Redis or a relational store.
  \item \textbf{Model retraining}: At 1{,}200 accounts, full retraining takes
    $\sim$60--90 seconds. As the dataset grows, consider incremental learning or
    dedicated ML infrastructure (MLflow + a training VM).
  \item \textbf{Monitoring}: Add logging to a structured sink (e.g., ELK stack
    or Azure Monitor) to track API latency, prediction drift, and error rates.
\end{itemize}

\subsection{Security Checklist}

\begin{itemize}[leftmargin=1.5cm, noitemsep]
  \item[$\square$] TLS/HTTPS enforced on all endpoints
  \item[$\square$] Authentication \& role-based access (Analyst / Executive / Admin)
  \item[$\square$] CORS restricted to known origins (remove wildcard \texttt{*})
  \item[$\square$] Input validation on \texttt{/api/simulate} payload
  \item[$\square$] CSV files stored outside web root or in private object storage
  \item[$\square$] Dependency audit (\texttt{pip-audit}) run before each release
\end{itemize}

% ═══════════════════════════════════════════════════════════════════════════════
\section{Appendix A — Data Dictionary}
% ═══════════════════════════════════════════════════════════════════════════════

\begin{longtable}{p{4.5cm}p{3.5cm}p{5cm}}
\toprule
\textbf{Field} & \textbf{Source} & \textbf{Description} \\
\midrule
\endhead
account\_external\_id & Account.csv & Primary key, e.g.\ \texttt{ACC-00001} \\
loyalty\_tier\_label  & Account.csv & Current CRM tier: Bronze / Silver / Gold \\
clv\_2025             & step\_02\_features.csv & Actual 2025 revenue (Kč) — training target \\
clv\_2025\_predicted  & Model output & XGBoost regression prediction (Kč) \\
churn\_probability    & step\_08\_churn\_output.csv & $P(\text{churn}) \in [0,1]$ \\
churn\_risk           & Derived & Low / Medium / High (thresholds: 35\%, 65\%) \\
suggested\_tier       & Derived & Bronze / Silver / Gold from predicted CLV \\
tier\_match           & Derived & Boolean — actual tier = suggested tier \\
recency\_days         & step\_02\_features.csv & Days since last purchase \\
frequency             & step\_02\_features.csv & Number of purchase events \\
monetary\_total       & step\_02\_features.csv & Total historical spend (Kč) \\
spend\_trend\_1y      & step\_02\_features.csv & YoY spend change (ratio) \\
email\_open\_rate     & Activity\_\_c.csv & Campaign email open rate (\%) \\
app\_usage\_score     & Activity\_\_c.csv & Normalised mobile app engagement score \\
\bottomrule
\caption{Key data dictionary (selected fields)}
\end{longtable}

% ═══════════════════════════════════════════════════════════════════════════════
\section{Appendix B — Segment Output Example}
% ═══════════════════════════════════════════════════════════════════════════════

\begin{table}[H]
\centering\small
\begin{tabular}{lrrrr}
\toprule
\textbf{Segment} & \textbf{Accounts} & \textbf{Avg CLV (Kč)} & \textbf{Share (\%)} \\
\midrule
Champions              & 58  & 62{,}410 & 4.8 \\
High-Value At-Risk     & 6   & 9{,}790  & 0.5 \\
Loyal Mid-Tier         & 287 & 11{,}240 & 23.9 \\
Dormant High-Potential & 213 & 8{,}520  & 17.8 \\
Growing New            & 134 & 7{,}830  & 11.2 \\
\bottomrule
\end{tabular}
\caption{Representative segment output (values approximate)}
\end{table}

% ─── End ──────────────────────────────────────────────────────────────────────
\vfill
\begin{center}
  \color{sfgray}\small
  CLV Prediction Platform — Technical Report v1.0 · \today \\
  Confidential — Internal Use Only
\end{center}

\end{document}
